\subsection{Mixed lower bounds}
\label{sec:nonpluripolar-common-measure-mixed}
The bounded common-measure inequality passes through the same finite-level restrictions that define the non-pluripolar product. A polar part of the reference measure causes no difficulty: domination by the pure non-pluripolar measures makes both prescribed densities vanish there.

\begin{proposition}\label{prop:nonpluripolar-common-measure-mixed}
Let $T,S$ be closed positive $(1,1)$-currents on compact K\"ahler $X$ of dimension $n\geq1$. Let $\nu$ be a finite positive Radon measure, and suppose $f,g\geq0$ belong to $L^1(\nu)$ with
\[
 \langle T^n\rangle\geq f\nu,\qquad
 \langle S^n\rangle\geq g\nu.
\]
For $1\leq k\leq n-1$,
\[
 \langle T^k\wedge S^{n-k}\rangle
                    \geq f^{k/n}g^{(n-k)/n}\nu.
\]
The endpoint inequalities are the hypotheses. All these products are finite by \cref{thm:compact-nonpluripolar-finiteness}.
\end{proposition}
\begin{proof}
Work on a relatively compact coordinate domain where $T=\ddc a$ and $S=\ddc b$ with psh potentials bounded above. For $j\geq1$ put
\[
 a_j=\max\{a,-j\},\quad b_j=\max\{b,-j\},\quad
 G_j=\{a>-j\}\cap\{b>-j\}.
\]
On $G_j$, truncation locality identifies each pure bounded product with the corresponding non-pluripolar measure. Positivity therefore gives on the whole coordinate domain
\[
 (\ddc a_j)^n\geq\mathbf1_{G_j}f\nu,\qquad
 (\ddc b_j)^n\geq\mathbf1_{G_j}g\nu.
\]
Apply \cref{thm:bounded-common-measure-mixed} to these bounded potentials and then restrict to $G_j$. The defining truncation identity gives
\[
 \langle T^k\wedge S^{n-k}\rangle
                  \geq\mathbf1_{G_j}f^{k/n}g^{(n-k)/n}\nu.
\]
The complement of $\bigcup_jG_j$ is the union of the two local pole sets. Both pure non-pluripolar measures give this union mass zero, since they annihilate every pluripolar set. The assumed domination consequently gives $f=g=0$ $\nu$-almost everywhere there. Monotone convergence on nonnegative compact tests now removes the indicator. Coordinate domains cover $X$, proving the assertion. For $n=1$ there is no interior mixed term.
\end{proof}
\index{non-pluripolar product!common-measure mixed inequality}

\begin{corollary}\label{cor:nonpluripolar-common-measure-convex}
Under these hypotheses, for $0\leq t\leq1$,
\[
 \left\langle\big((1-t)T+tS\big)^n\right\rangle
       \geq\big((1-t)f^{1/n}+tg^{1/n}\big)^n\nu.
\]
\end{corollary}
\begin{proof}
Positive multilinearity expands the left side. Apply the proposition to the interior mixed coefficients and use the two pure bounds for the endpoints. The binomial formula gives the right side. Endpoint values of $t$ use the pure hypotheses directly; for $n=1$ use additivity. Weighted arithmetic--geometric mean bounds each interior density by a multiple of $f+g$, so the right-hand measure is finite.
\end{proof}
